\chapter{Бинарно-тетраэдральный carrier gate}

Внешний exploratory-аудит предложил заменить
\(\mathbb{RP}^3=S^3/\mathbb Z_2\) на сферическую форму
\begin{equation}
 M_{2T}=S^3/2T,
 \qquad
 2T\simeq\operatorname{SL}(2,3),
 \qquad |2T|=24,
\end{equation}
и интерпретировать число \(23\) как размерность augmentation ideal
\(\mathbb C[2T]\ominus\mathbf1\). Проверим отдельно carrier-нормировку,
спектр и представленческий мост.

\section{Carrier-тест}

При единичном радиусе
\begin{equation}
 \Vol(M_{2T})=\frac{2\pi^2}{24}=\frac{\pi^2}{12}.
\end{equation}
Тем самым исходное условие \(Z_A=\pi^2\) нарушено в двенадцать раз.
Абелианизация \(2T\) равна \(\mathbb Z_3\), поэтому плоские абелевы
характеры имеют фазовый шаг \(2\pi/3\), а не \(\pi\). Кроме того, минимальный
геометрический сдвиг задаётся элементом с вещественной частью \(1/2\), так
что систола круглого quotient-а равна
\begin{equation}
 \ell_{sys}(M_{2T})=\arccos(1/2)=\frac{\pi}{3}.
\end{equation}
Прямая замена в прежнем геометрическом скаляре даёт
\begin{align}
 S_{geo}^{2T}
 &=4\pi^3+\frac{\pi^2}{12}+\frac{\pi}{3}
 =125.8947713058\ldots,\\
 S_{vac}^{2T}
 &=S_{geo}^{2T}-\frac{1}{24S_{geo}^{2T}}
 -\frac{1}{\pi^4(S_{geo}^{2T})^2}
 =125.8944396939\ldots.
\end{align}
Расхождение с train-якорем \(\alpha^{-1}\) равно
\begin{equation}
 S_{vac}^{2T}-\alpha^{-1}=-11.1415594831\ldots.
\end{equation}
Следовательно, candidate не является продолжением исходной numerical
architecture без нового множителя или нового определения \(S_{geo}\).

\section{Явное усреднение характеров}

Для \(g\in2T\subset SU(2)\) характер
\(\operatorname{Sym}^{\ell}(\mathbb C^2)\) равен
\(U_\ell(\Re g)\), где \(U_\ell\) --- полином Чебышёва второго рода.
Поэтому
\begin{equation}
 m_\ell^{2T}
 =\dim\operatorname{Sym}^{\ell}(\mathbb C^2)^{2T}
 =\frac1{24}\sum_{g\in2T}U_\ell(\Re g).
\end{equation}
Прямой подсчёт по двадцати четырём quaternion units даёт первые ненулевые
степени
\begin{equation}
 \ell=0,6,8,12,14,16,18,\ldots.
\end{equation}
Первая положительная scalar shell имеет
\begin{equation}
 \lambda_6=6(6+2)=48,
 \qquad d_6=(6+1)m_6^{2T}=7.
\end{equation}

Для coexact one-forms homogeneous left quotient multiplicity равна
\begin{equation}
 d_n^{2T}
 =n\,m_{n+1}^{2T}+(n+2)m_{n-1}^{2T}.
\end{equation}
Первый уровень сохраняется:
\begin{equation}
 n=1,
 \qquad \lambda_1=4,
 \qquad d_1^{2T}=3.
\end{equation}
Поэтому положительный Bessel-tail не исчезает. При \(R_3=R_1=1\)
получаем
\begin{equation}
 T_{coex}^{2T}=7.6133956135\cdot10^{-6}>0,
\end{equation}
причём первый уровень даёт более
\(0.9999999999\) полной суммы. Замена quotient-а лишь примерно делит
доминирующий residual пополам и не создаёт cancellation theorem.

\section{Почему augmentation ideal не является quotient spectrum}

Регулярное представление действительно имеет каноническое разложение
\begin{equation}
 \mathbb C[2T]=\mathbf1\oplus I_{aug},
 \qquad \dim I_{aug}=23.
\end{equation}
Однако функции и формы на \(S^3/2T\) соответствуют \(2T\)-инвариантным
состояниям на \(S^3\). Quotient-проектор сохраняет \(\mathbf1\)-isotypic
часть и удаляет nontrivial regular components. Следовательно,
\(I_{aug}\) не возникает как конечный physical spectral sector самого
quotient-а.

Утверждение, что \(Q_8\) не действует на этом пространстве, также неверно:
\(Q_8\subset2T\) действует левыми сдвигами. Ограничение регулярного
представления содержит три копии регулярного \(Q_8\)-модуля, поэтому
\begin{equation}
 \dim\mathbb C[2T]^{Q_8}=3,
 \qquad
 \dim I_{aug}^{Q_8}=2.
\end{equation}
Старый spinor-line no-go здесь непосредственно не применяется, но это не
создаёт требуемого spectral bridge.

\begin{nogocriterion}[Binary-tetrahedral direct-carrier no-go]
Прямая замена \(\mathbb{RP}^3\) на \(S^3/2T\) нарушает locked volume и
flat-phase tests, сохраняет положительный coexact tail и не реализует
augmentation ideal как quotient eigenspace. Поэтому маршрут закрыт как
геометрический источник физического ранга \(23\). Он может быть переоткрыт
только как новая модель с отдельно выведенным внутренним
\(\mathbb C[2T]\)-fiber и единым действием, связывающим его с Dirac,
калибровочным и normalization-sensitive секторами.
\end{nogocriterion}

Машинный аудит сохранён в
\texttt{s2t\_binary\_tetrahedral\_carrier\_results.json}.